\section{GPU 后端优化细讲：device buffer、stream、kernel launch 与端到端延迟}

\subsection{GPU 后端不是把 CPU 指针换成 device pointer}

GPU 对 AI 推理很重要，但它也很容易被误解。初学者常以为只要把权重放到显存、把算子换成 GPU kernel，就一定更快。真实情况更细：大 prefill 可能非常适合 GPU，因为矩阵乘规模大、并行度高；decode 单 token 却可能被 kernel launch、同步、KV cache 读流、logits 回传和 sampler 边界限制。GPU 峰值 FLOPS 很高，不等于一次本地对话的端到端 latency 一定低。

从编译器角度看，GPU 后端是另一个 lowering 目标。它和 CPU 后端共享模型语义、shape、dtype、quant descriptor、KV cache 合同和 oracle，但物理执行完全不同。CPU 后端关心 cache line、SIMD 寄存器和线程池；GPU 后端关心 device memory、coalesced load、shared memory、warp/block、stream、kernel launch 和 host-device staging。

\begin{keyidea}
GPU 后端优化的核心是端到端计划，不是单个 kernel 时间。device buffer、stream、workspace、staging、同步点和 fallback 都必须进入 executable plan 和 profiler。
\end{keyidea}

本节不会依赖某个具体 GPU API。无论使用 CUDA、HIP、Metal、Vulkan compute 还是厂商库，工程边界都类似：资源要 RAII 管理，host/device 所有权要明确，异步执行要有 event 证据，kernel launch 要进入 profiler，fallback 不能悄悄打断语义合同。

\subsection{GPU 后端的 C++20 RAII 边界}

GPU 后端首先要把资源生命周期写清楚。下面这些对象不应该以裸句柄和裸指针散落在 runtime 中：

\begin{lstlisting}[numbers=none,caption={GPU 后端 RAII 对象}]
GpuContext:
  owns device selection and backend capability query

DeviceBuffer:
  owns device memory allocation and free

Stream:
  owns asynchronous execution queue

Event:
  owns timestamp/synchronization marker

Workspace:
  owns temporary device memory for algorithms

StagingBuffer:
  owns pinned or mapped host memory for transfers

KernelAdapter:
  binds executable plan descriptors to launch calls
\end{lstlisting}

这些对象的价值不是包装 API，而是建立工程不变量。DeviceBuffer 的析构释放显存；Stream 的生命周期覆盖异步 kernel；Event 让 profiler 能测真实 device 时间；StagingBuffer 显式表示 host-device 边界；KernelAdapter 把高层 descriptor 转换成具体 launch 参数。没有这些边界，GPU 后端很快会变成一堆到处传递的 \code{void*} 和句柄。

\begin{warningbox}
异步后端最怕生命周期伪装。CPU 代码返回不代表 GPU kernel 已经读完输入；arena 或 staging buffer 不能在相关 stream/event 完成前复用。
\end{warningbox}

\subsection{host/device 驻留：数据在哪里比函数在哪里更重要}

GPU 推理的第一条账本是数据驻留。权重、packed weight、activation、KV cache、logits、token id、sampler state 分别在 host 还是 device？什么时候传输？传输是否同步？这些问题往往比单个 kernel 代码更影响端到端延迟。

\begin{center}
\begin{tabular}{p{0.18\linewidth}p{0.26\linewidth}p{0.36\linewidth}}
\toprule
数据 & 理想驻留 & 风险 \\
\midrule
packed weight & device 常驻 & 首次上传/显存峰值/多 GPU 分片 \\
activation & device arena & 跨 stream 生命周期和 workspace 冲突 \\
KV cache & device 常驻或分层 & 长上下文显存占用、page table、量化 \\
logits & device 计算，少量回传 & 完整 vocab 回传太贵，sampler 边界同步 \\
token id & host/device 双方可见 & 每步小传输和同步放大 latency \\
\bottomrule
\end{tabular}
\end{center}

一个常见坏路径是：每个 decode step 在 GPU 上算 logits，然后把完整 vocab logits 拷回 CPU，再在 CPU 上做 sampler，再把下一个 token 拷回 GPU。对于大词表，这个回传可能成为固定开销。更好的计划可能是在 GPU 上做 top-k 或 sampler 的一部分，只回传 next token 或少量调试数据。但这会改变 sampler 的执行位置，必须进入 executable plan 和 oracle。

\subsection{stream 和 event：异步不是自动并行}

GPU API 往往是异步的。向 stream 提交 kernel 后，CPU 线程可以继续走；但这不等于工作已经完成，也不等于不同操作可以随便重排。stream 和 event 是计划中的同步语义。

\begin{lstlisting}[numbers=none,caption={GPU plan 中需要记录的同步边界}]
GpuPlan:
  stream decode_stream
  events:
    after_kv_append
    after_attention
    before_logits_staging
  dependencies:
    attention waits for kv_append
    sampler waits for logits staging
    arena slot reuse waits for producing kernels
\end{lstlisting}

profiler 也必须区分 CPU 提交时间和 GPU 执行时间。只测 kernel 内部时间，会漏掉 launch overhead、host-device copy、stream wait、CPU sampler 和计划选择。只测 CPU 墙钟时间，又难以定位慢在 device kernel 还是同步。两者都要记录。

\begin{systemdiagram}{GPU executable plan 中的 stream、event 与 staging}
\begin{pdfdiagram}[x=1cm,y=1cm]
  \node[book accent node,text width=2.25cm] (host) at (0,2.15) {host runtime\\select plan};
  \node[book teal node,text width=2.25cm] (launch) at (3.0,2.15) {launch kernels\\stream};
  \node[book purple node,text width=2.25cm] (dev) at (6.0,2.15) {device compute\\attention/MLP};
  \node[book amber node,text width=2.25cm] (stage) at (9.0,2.15) {staging\\logits/top-k};

  \node[book navy node,text width=2.25cm] (event) at (6.0,0.35) {events\\timestamps};
  \node[book red node,text width=2.25cm] (sample) at (9.0,0.35) {sampler\\sync boundary};

  \draw[book strong arrow] (host) -- (launch);
  \draw[book strong arrow] (launch) -- (dev);
  \draw[book strong arrow] (dev) -- (stage);
  \draw[book weak arrow] (dev) -- (event);
  \draw[book strong arrow] (stage) -- (sample);
  \draw[book weak arrow] (event) -- (sample);
\end{pdfdiagram}
\diagramnote{本图要证明的命题是：GPU 后端的端到端延迟由 launch、device compute、event、staging 和 sampler 同时决定。}
\end{systemdiagram}

\subsection{kernel launch：小算子会被固定开销吞掉}

GPU kernel launch 有固定成本。prefill 的大 GEMM 能摊薄 launch 成本；decode 的许多小算子不一定能。如果每层都单独 launch RMSNorm、QKV、RoPE、KV append、attention、Wo、MLP gate、activation、down projection，decode 每个 token 会产生大量 launch 和同步。

解决方向有几类：

\begin{itemize}
  \item 算子融合：把 RMSNorm 与 linear、RoPE 与 KV append 等合并。
  \item segment 调度：减少 host 侧逐算子解释和 profiler 细粒度开销。
  \item persistent 或批处理策略：对服务端多请求场景减少重复 launch。
  \item GPU sampler：减少 logits 完整回传和 host 同步。
\end{itemize}

融合仍然要受语义约束。比如 oracle 需要某层中间输出时，debug plan 不能完全隐藏物化点；量化权重需要 scale metadata，融合 kernel 必须正确读取；KV append 写入状态，后续 attention 必须看到正确 position。GPU fusion 不是“把越多函数塞进一个 kernel 越好”，而是减少内存流和 launch 开销，同时控制寄存器、shared memory 和调试边界。

\subsection{prefill GPU plan：吞吐和 workspace}

prefill 的输入 token 多，矩阵乘和 attention 更大，GPU 通常更有优势。编译器 lowering 要为 prefill 选择合适的 batched GEMM、attention kernel 和 workspace。workspace 是后端算法的临时内存，不能和 activation arena 随意混用。

\begin{lstlisting}[numbers=none,caption={prefill GPU plan 的主要字段}]
PrefillGpuPlan:
  prompt_length_bucket
  device_weight_layout
  activation_arena_size
  attention_workspace_size
  gemm_algorithm_id
  attention_algorithm_id
  stream_dependencies
  kv_cache_write_layout
\end{lstlisting}

prefill 指标也要拆开。time to first token 包含 tokenizer、权重准备、KV cache 初始化、prefill graph 和 logits/sampler；prompt tokens/s 只描述 prompt 处理吞吐。若报告只写一个“GPU prefill 很快”，但没有说明 prompt 长度、batch、workspace、是否包含权重上传，就无法对标。

\subsection{decode GPU plan：小 batch、KV cache 与同步}

decode 的输入通常是一个 token。即使 GPU 算力强，单 token 路径也可能被固定开销限制。decode GPU plan 要重点处理三件事：

\begin{itemize}
  \item 减少每 token launch 和同步次数。
  \item 让 KV cache 读取尽量连续或分块，避免长上下文下显存读流退化。
  \item 减少 logits 回传，只同步 sampler 真正需要的数据。
\end{itemize}

\begin{lstlisting}[numbers=none,caption={decode GPU profiler 应拆分的事件}]
decode_step:
  host_plan_select_ns
  launch_overhead_ns
  rmsnorm_linear_kernel_ns
  kv_append_kernel_ns
  attention_kernel_ns
  mlp_kernel_ns
  logits_staging_ns
  sampler_ns
  total_step_latency_ns
\end{lstlisting}

如果只看 \code{attention_kernel_ns}，可能会误以为 GPU 很快；如果 \code{logits_staging_ns + sampler_ns} 很大，用户感受到的 tokens/s 仍然不高。端到端报告必须包含 \code{total_step_latency_ns}。

\subsection{一个 decode step 的端到端时序}

GPU decode 不能只看 kernel profiler。单 token 生成是一条 host 和 device 共同参与的流水线。即使每个 device kernel 都很快，host 侧计划选择、kernel launch、event wait、logits staging 和 sampler 也可能支配端到端延迟。

\begin{lstlisting}[numbers=none,caption={GPU decode step 的最小时序}]
host:
  bind current token and position
  submit rmsnorm/qkv/rope/attention/mlp kernels
  submit logits or top-k staging
  wait for sampler boundary
  sample next token

device stream:
  rmsnorm_qkv_kernel
  rope_kv_append_kernel
  attention_decode_kernel
  mlp_kernel
  logits_topk_kernel
  copy top-k or logits slice to host
\end{lstlisting}

这条时序里至少有三个同步边界。第一，KV append 与 attention 之间要有顺序，attention 必须看到当前 token 写入后的 K/V。第二，logits 或 top-k 结果回到 host 前，sampler 不能读。第三，arena 或 workspace 复用前，相关 stream 上的 kernel 必须完成。若这些边界只靠“代码看起来按顺序写”维持，就很容易在异步后端出错。

\begin{lstlisting}[numbers=none,caption={StepTrace 应记录的字段}]
StepTrace:
  request_id
  step_index
  model_position
  context_length
  selected_plan_id
  host_submit_begin_ns
  host_submit_end_ns
  device_kernel_begin_ns
  device_kernel_end_ns
  staging_begin_ns
  staging_end_ns
  sampler_begin_ns
  sampler_end_ns
  sync_wait_ns
  total_step_latency_ns
\end{lstlisting}

有了这个结构，报告才能回答具体问题：慢在 host 提交、device kernel、staging、wait，还是 sampler。没有这个拆分，只拿 GPU 工具里某个 kernel 的耗时去推断用户感受到的 tokens/s，是不完整的。

\subsection{GPU decode 的最小可用优化顺序}

GPU 后端优化也要分阶段，不要一开始就写一个巨大的 fused kernel。一个更稳的顺序是：

\begin{enumerate}
  \item 先让所有权重和 KV cache 常驻 device，消除每步大块 host-device 传输。
  \item 把 logits staging 从完整 vocab 缩小到 sampler 需要的最小集合，或把 top-k 放到 device。
  \item 合并小算子，优先融合 RMSNorm+Linear、RoPE+KV append 这类语义边界清晰的路径。
  \item 对 attention decode 按 context length 分档，短上下文、长上下文、paged KV 使用不同 plan。
  \item 最后再做更激进的 persistent kernel、跨层 segment 或多请求 batch decode。
\end{enumerate}

每一级都要同时看 oracle 和 StepTrace。比如把 top-k 放到 GPU 后，logits 全量 oracle 可能不再自然可得，这时 debug plan 要保留可选的 full logits materialization；生产 plan 可以只回传 top-k，但必须声明 sampler 语义没有改变。

\subsection{GPU 量化：格式支持比位宽口号更重要}

GPU 上的低比特量化并不天然快。关键是目标硬件和库是否支持该格式，以及支持发生在哪一层。如果 int4 权重格式能直接进入高效矩阵乘路径，收益可能明显；如果每次都要先用自定义 kernel 解包成 fp16，再调用 GEMM，收益要重新计算。scale metadata 的布局也会影响 coalesced load。

\begin{center}
\begin{tabular}{p{0.22\linewidth}p{0.30\linewidth}p{0.32\linewidth}}
\toprule
路线 & 优点 & 风险 \\
\midrule
vendor quant kernel & 利用硬件和库优化 & 格式受限，descriptor 适配复杂 \\
自定义 dequant + GEMM & 实现直接，易验证 & 解包和额外写回可能很贵 \\
fused dequant matmul & 减少中间写回 & kernel 复杂，寄存器/shared memory 压力 \\
KV cache 量化 & 降低长上下文显存和带宽 & scale metadata 和 attention 精度风险 \\
\bottomrule
\end{tabular}
\end{center}

量化 verifier 要知道 GPU 后端实际支持什么。不能因为 manifest 写了 \code{q4}，就假设 GPU plan 可用。若 GPU 不支持该格式，应明确 fallback 到 CPU、fallback 到 dequant path，或拒绝该 plan，并在报告里写出原因。

\subsection{fallback：不能悄悄改变执行位置}

真实工程中，某些算子或格式 GPU 后端暂时不支持，需要 fallback。fallback 可以是 CPU reference、CPU optimized path、GPU 上的慢路径，或者 vendor 库之外的自定义 kernel。关键是不能悄悄发生。

\begin{lstlisting}[numbers=none,caption={fallback report 字段}]
FallbackEvent:
  operator_id
  reason: unsupported_dtype | unsupported_shape | workspace_limit | debug_oracle
  from_backend
  to_backend
  data_transfer_bytes
  added_sync_points
  correctness_profile
\end{lstlisting}

如果某层 attention fallback 到 CPU，会发生 device-to-host 拷贝、CPU 计算、host-to-device 拷贝和同步，端到端 latency 可能大幅上升。profiler 必须把它显示出来。oracle 也要知道 fallback 路径的误差合同可能不同。

\subsection{多 GPU 和分片：先建立单设备合同}

多 GPU 推理涉及权重分片、KV cache 分布、通信、pipeline/tensor parallel 和跨设备同步。本册后续可以扩展，但工程顺序应该先把单设备合同建立好。原因很简单：如果单 GPU 的 descriptor、stream、workspace、oracle 和 profiler 都不清楚，多 GPU 只会把错误放大。

单设备稳定后，多 GPU 至少要新增这些事实：

\begin{itemize}
  \item 张量分片轴和每个 shard 的 shape。
  \item 跨设备通信算子和同步点。
  \item KV cache 是复制、分片还是按层分布。
  \item logits 聚合和 sampler 所在设备。
  \item profiler 能拆分 compute、communication 和 wait。
\end{itemize}

这些内容和第二册分布式边界、背压、运行时调度有关。它们不是本节主要目标，但要从一开始为 descriptor 和 plan 留出扩展空间。

\subsection{GPU 后端验收报告}

GPU 后端优化的报告至少要覆盖：

\begin{lstlisting}[numbers=none,caption={GPU 后端验收报告字段}]
hardware:
  gpu_model
  driver_runtime_version
  device_memory
  enabled_backend_features

plan:
  backend_plan_id
  device_layouts
  stream_count
  workspace_bytes
  fallback_events

performance:
  weight_upload_time_ms
  pack_or_convert_time_ms
  prefill_tokens_per_s
  time_to_first_token_ms
  decode_tokens_per_s
  total_step_latency_ms
  launch_time_ms
  staging_time_ms
  device_kernel_breakdown

correctness:
  oracle_profile
  logits_error
  topk_changed
  generated_token_mismatch
\end{lstlisting}

注意 \code{launch_time_ms} 和 \code{staging_time_ms} 必须出现。它们正是很多 GPU decode 路径慢的原因。如果报告只有 device kernel breakdown，就还不是端到端推理报告。

\subsection{本节验收线}

读完本节，应该能说明 GPU 后端的真实工程边界：

\begin{itemize}
  \item GPU 后端共享 AI 编译器的语义合同，但拥有自己的 device buffer、stream、event、workspace、staging 和 kernel adapter。
  \item 数据驻留和 host-device 传输常常比单个 kernel 更影响端到端 latency。
  \item stream/event 必须进入 executable plan，arena 和 staging buffer 复用要等待异步工作完成。
  \item kernel launch 是 decode 小算子的固定成本，fusion 和 segment 调度要围绕它设计。
  \item prefill 关注吞吐和 workspace，decode 关注 launch、KV cache 读流、logits staging 和 sampler 同步。
  \item GPU 量化要看硬件和库是否真正支持目标格式，fallback 必须显式报告。
  \item GPU benchmark 必须同时报告 device kernel 时间、launch、staging、同步、端到端 tokens/s 和 oracle。
\end{itemize}

到这里，CPU 和 GPU 后端有了共同的工程语言：descriptor 表达语义，lowering 选择物理计划，runtime 绑定状态，oracle 证明正确，profiler 证明性能。两者的差异不是谁更“高级”，而是机器账本不同。
